\subsection{Безразмерные параметры микрорельефа}

Для параметрического анализа удобно перейти к безразмерным параметрам текстуры.
Используются масштабы, введённые в подразделе~2.1.6
(таблица~\ref{tab:nondim_scales}): характерная толщина зазора~$h_0$
и характерные длины~$L_1$,~$L_2$ по поверхностным координатам~$x_1$,~$x_2$.
Безразмерные координаты и уравнение Рейнольдса уже заданы
в подразделе~2.1.6; здесь нормируются только параметры текстуры.

Ключевым безразмерным параметром является отношение глубины элемента
к характерной толщине зазора:
\begin{equation}\label{eq:Hp_nondim}
  H_p = \frac{h_p}{h_0}.
\end{equation}

Требование применимости уравнения Рейнольдса к текстурированной
поверхности (подраздел~2.1.5) формулируется как $H_p \ll 1$.

Безразмерные полуоси элемента определяются делением на соответствующие
масштабы длины:
\begin{equation}\label{eq:AB_nondim}
  A^{*} = \frac{a}{L_2}, \quad B^{*} = \frac{b}{L_1},
\end{equation}
аспектное отношение элемента:
\begin{equation}\label{eq:kappa}
  \kappa = \frac{a}{b}.
\end{equation}

При $\kappa = 1$ элемент осесимметричен (круговая лунка);
при $\kappa \neq 1$ -- эллиптический.
Заметим, что при $L_1 \neq L_2$ выполняется
\[
  \kappa = \frac{a}{b}
         = \frac{A^{*}}{B^{*}}\,\frac{L_2}{L_1},
\]
а при $L_1 = L_2$ получаем $\kappa = A^{*}/B^{*}$.

Безразмерные шаги раскладки:
\begin{equation}\label{eq:P_nondim}
  P_1^{*} = \frac{p_1}{L_1}, \quad P_2^{*} = \frac{p_2}{L_2}.
\end{equation}

Условия непересечения элементов (подраздел~2.2.1) удобно проверять
через отношения $p_1/(2b)$ и $p_2/(2a)$, которые показывают,
во сколько раз шаг превышает полный размер элемента
($2b$ по $x_1$ и $2a$ по $x_2$).

Коэффициент заполнения~$\phi$ (подраздел~2.2.1) является безразмерным.
Для эллиптических элементов на регулярной решётке при непересекающихся
элементах и пренебрежении краевыми эффектами
\begin{equation}\label{eq:phi_nondim}
  \phi \approx \frac{\pi\, a\, b}{p_1\, p_2}
       = \frac{\pi\, A^{*}\, B^{*}}{P_1^{*}\, P_2^{*}}.
\end{equation}

Полный перечень безразмерных параметров микрорельефа сведён в табл.~\ref{tab:nondim_texture}.

Ряд параметров текстуры является безразмерным по определению:
$\rho_0 \in (0,\,1)$ -- параметр плоского дна (подраздел~2.2.2);
$\alpha$ -- угол ориентации элемента.
Для канавок вводится безразмерная длина $\ell^{*} = \ell / L_1$;
для сквозных канавок по периодическому направлению при выборе
$L_1$ равным длине периода $\ell^{*} = 1$.
Для повёрнутых элементов ($\alpha \neq 0$) используются локальные
координаты из подраздела~2.2.2; переход к безразмерным величинам
выполняется делением длин на соответствующие масштабы из~2.1.6.

\begin{table}[!ht]
\centering
\caption{Безразмерные параметры микрорельефа}
\label{tab:nondim_texture}
\begin{tabularx}{\textwidth}{|X|c|c|}
\hline
\multicolumn{1}{|c|}{\textbf{Параметр}} & \textbf{Обозначение} & \textbf{Определение} \\
\hline
Безразмерная глубина             & $H_p$    & $h_p / h_0$                          \\
\hline
Безразмерная полуось по $x_2$    & $A^{*}$  & $a / L_2$                            \\
\hline
Безразмерная полуось по $x_1$    & $B^{*}$  & $b / L_1$                            \\
\hline
Отношение полуосей               & $\kappa$ & $a / b$                              \\
\hline
Безразмерный шаг по $x_1$        & $P_1^{*}$ & $p_1 / L_1$                        \\
\hline
Безразмерный шаг по $x_2$        & $P_2^{*}$ & $p_2 / L_2$                        \\
\hline
Коэффициент заполнения           & $\phi$   & $N A_{\mathrm{elem}} / A_{\mathrm{tex}}$ \\
\hline
Параметр плоского дна            & $\rho_0$ & $(0,\,1)$                            \\
\hline
Угол ориентации                  & $\alpha$ & --                                   \\
\hline
Безразмерная длина канавки       & $\ell^{*}$ & $\ell / L_1$                       \\
\hline
\end{tabularx}
\end{table}

\FloatBarrier

Введённый набор безразмерных параметров
($H_p$, $A^{*}$, $B^{*}$, $\kappa$, $P_1^{*}$, $P_2^{*}$, $\phi$, $\rho_0$, $\ell^{*}$)
позволяет проводить параметрические исследования и сравнивать результаты
для различных типов подшипников при одинаковых безразмерных группах.

Физический смысл каждого безразмерного параметра элемента
микрорельефа целесообразно кратко прокомментировать.
Глубина~$H_p = h_p/h_0$ задаёт масштаб добавки к зазору
относительно характерной толщины плёнки. При $H_p \ll 1$
лунка неглубокая; её гидродинамический эффект мал, так как
создаваемый ею дополнительный клин пропорционален $H_p$.
При $H_p \gtrsim 1$ глубина лунки сопоставима с зазором,
что может нарушить течение в соседних ячейках и снизить
несущую способность; для эллипсоидальных лунок оптимальная
глубина, как правило, находится в диапазоне
$H_p \approx 0{,}2$--$0{,}5$.

Безразмерные полуоси~$A^{*}$, $B^{*}$ задают размеры элемента
относительно длин расчётной области в каждом из координатных
направлений. Произведение~$\pi A^{*} B^{*}$ пропорционально
безразмерной площади одной лунки; коэффициент
покрытия~$\phi$ определяет суммарную долю текстурированной
площади. Ориентация эллипса ($A^{*} \neq B^{*}$) влияет
на анизотропию гидродинамического эффекта и, следовательно,
на соотношение между несущей способностью, трением и утечками.

Шаги~$P_1^{*}$, $P_2^{*}$ (периоды раскладки по~$s_1$
и~$s_2$ соответственно) совместно с~$A^{*}$, $B^{*}$
и~$\phi$ образуют замкнутый набор параметров текстуры,
достаточный для однозначного задания шахматной раскладки.
При шахматной схеме реальный шаг в нечётных рядах сдвигается
на $P_1^{*}/2$, что обеспечивает более равномерное перекрытие
поверхности при той же плотности.

\textit{Конкуренция двух гидродинамических эффектов.}
Влияние микроструктур на~несущую способность определяется конкуренцией двух противоположно направленных механизмов. Первый механизм~--- генерация давления: передняя стенка лунки создаёт дополнительный сходящийся клин, что усиливает гидродинамический эффект и~повышает давление над~лункой. Второй механизм~--- сброс давления: задняя стенка лунки образует расходящийся клин, где давление падает; при~глубоких лунках в~этой зоне формируется локальная кавитационная область, которая обрывает возможность дальнейшего нарастания давления. Суммарный эффект на~несущую способность определяется соотношением этих двух вкладов.

При~малых глубинах~$H_p \ll 1$ первый механизм преобладает, и~несущая способность монотонно возрастает с~ростом~$H_p$. При~чрезмерно глубоких лунках~$H_p \gtrsim 1$ кавитационная зона занимает значительную часть задней стенки, суммарный вклад становится отрицательным, и~несущая способность начинает снижаться относительно гладкой поверхности. Следствием является существование оптимальной глубины~$H_p^{*}$, при~которой несущая способность максимальна. Для~эллипсоидальных лунок в~типичных конфигурациях этот оптимум находится в~диапазоне $H_p^{*} \approx 0{,}2$--$0{,}5$~--- именно эти значения использованы в~главах~5--7.

Аналогичная конкуренция наблюдается по~параметру коэффициента покрытия~$\phi$. При~малых~$\phi$ лунки редки и~суммарный гидродинамический эффект мал; при~$\phi \to 1$ лунки занимают практически всю поверхность, клиновой эффект между лунками исчезает и~несущая способность снова снижается. Оптимальный коэффициент покрытия для~радиальных и~упорных подшипников, по~данным численных и~экспериментальных исследований, как~правило, находится в~диапазоне $\phi \approx 0{,}1$--$0{,}4$, хотя точное значение зависит от~типа подшипника, режима нагружения и~формы элементов микрорельефа.

\textit{Конкуренция двух гидродинамических эффектов.}
Влияние микроструктур на~несущую способность определяется конкуренцией двух противоположно направленных механизмов. Первый механизм --- генерация давления: передняя стенка лунки создаёт дополнительный сходящийся клин, что усиливает гидродинамический эффект и~повышает давление над~лункой. Второй механизм --- сброс давления: задняя стенка лунки образует расходящийся клин, где давление падает; при~глубоких лунках в~этой зоне формируется локальная кавитационная область, которая обрывает возможность дальнейшего нарастания давления. Суммарный эффект на~несущую способность определяется соотношением этих двух вкладов.

При~малых глубинах~$H_p \ll 1$ первый механизм преобладает, и~несущая способность монотонно возрастает с~ростом~$H_p$. При~чрезмерно глубоких лунках~$H_p \gtrsim 1$ кавитационная зона занимает значительную часть задней стенки, суммарный вклад становится отрицательным, и~несущая способность начинает снижаться относительно гладкой поверхности. Следствием является существование оптимальной глубины~$H_p^{*}$, при~которой несущая способность максимальна. Для~эллипсоидальных лунок в~типичных конфигурациях этот оптимум находится в~диапазоне $H_p^{*} \approx 0{,}2$--$0{,}5$~--- именно эти значения использованы в~главах~5--7.

Аналогичная конкуренция наблюдается по~параметру коэффициента покрытия~$\phi$. При~малых~$\phi$ лунки редки и~суммарный гидродинамический эффект мал; при~$\phi \to 1$ лунки занимают практически всю поверхность, клиновой эффект между лунками исчезает и~несущая способность снова снижается. Оптимальный коэффициент покрытия для~радиальных и~упорных подшипников, по~данным численных и~экспериментальных исследований, находится в~диапазоне $\phi \approx 0{,}1$--$0{,}4$, хотя точное значение зависит от~типа подшипника, режима нагружения и~формы элементов микрорельефа~\cite{senatore2018}.
